\documentclass{report}

\usepackage[latin1]{inputenc}
\usepackage[T1]{fontenc}
\usepackage[francais]{babel}
\usepackage[top=2.5cm, bottom=2.5cm, left=2.5cm, right=2.5cm]{geometry}
\usepackage{graphicx}
\usepackage{amsmath,amsfonts,amssymb}
\usepackage{parallel}
\usepackage{float}
\restylefloat{table}

\usepackage[linesnumbered,algoruled,boxed,lined]{algorithm2e}

\setlength\parindent{24pt} % indentation paragraphes

\begin{document}
\begin{titlepage}

\newcommand{\HRule}{\rule{\linewidth}{0.5mm}} % Defines a new command for the horizontal lines, change thickness here

\center % Center everything on the page
 
%----------------------------------------------------------------------------------------
%	HEADING SECTIONS
%----------------------------------------------------------------------------------------

\textsc{\LARGE ENSEEIHT}\\[1.5cm] % Name of your university/college
\textsc{\Large Sémantique et TDL}\\[0.5cm] % Major heading such as course name
\textsc{\large Projet}\\[0.5cm] % Minor heading such as course title

%----------------------------------------------------------------------------------------
%	TITLE SECTION
%----------------------------------------------------------------------------------------

\HRule \\[0.4cm]
{ \huge \bfseries  Compilateur du langage $\mu$C^{\#}}\\[0.4cm] % Title of your document
\HRule \\[1.5cm]
 
%----------------------------------------------------------------------------------------
%	AUTHOR SECTION
%----------------------------------------------------------------------------------------

\emph{auteurs:}\\
\textsc{achard} - \textsc{vincent} \\
\textsc{barras} - \textsc{thomas} \\
\textsc{bouqsimi} - \textsc{amine} \\
\textsc{rigondaud} - \textsc{clément} \\

%----------------------------------------------------------------------------------------
%	DATE SECTION
%----------------------------------------------------------------------------------------

\date{} %enleve la date

%----------------------------------------------------------------------------------------
%	LOGO SECTION
%----------------------------------------------------------------------------------------

\vspace{7.2cm}
\begin{center}
%\includegraphics[width=5]{N7.png}
\end{center}

%----------------------------------------------------------------------------------------

\vfill % Fill the rest of the page with whitespace

\end{titlepage}

\tableofcontents %table des matières.

\vfill

\part{Langage $\mu$C}

\chapter{Choix de conception}

\begin{itemize}
	\item Nous avons choisi de gérer le problème d'insertion de doublons directement dans la TDS afin de rendre plus clair le contenu des actions sémantiques.
	\item La TDS d'un bloc est créée à l'"extérieur" de celui-ci (1\iere{} apparition du non-terminal BLOC).
	\item Nous avons choisi de rajouter des méthodes à la machine TAM, celles-ci permettant de faciliter la gestion du déplacement dans le bloc courant.
	\item La classe INFOFONC a été créée afin de faciliter la manipulation des fonctions: paramètres et retour.
	\item Nous stockons les paramètres d'une fonction à la déclaration dans une unique TDS pour relever les erreurs de paramètres qui ont le même nom lors de l'appel d'une fonction.
	\item de nouvelles fonctions compareTo ont été définies afin de simplifier la vérification de type entre opérateurs et types.  
\end{itemize}

\chapter{Limitations}
\begin{itemize}
	\item La génération de code n'est pas terminée, nous nous sommes focalisés sur la bonne gestion de la TDS et des types.
	\item Des exemples simples de code sont en cours de création et de traduction en assembleur, pour simplifier le travail de vérification de génération de code dans le futur.
	
\end{itemize}



\end{document}